\documentclass{beamer}

\usetheme{metropolis}

\title{Miller-Rabin Primality Test}
\author{Siddharth \& Amit}
\date{\today}

\newcommand{\Z}{\mathbb{Z}}
\newcommand{\bld}[1]{\textbf{#1}}
\newcommand{\itl}[1]{\textit{#1}}
\newcommand{\N}{\mathbb{N}}
\newcommand{\bigO}{\mathcal{O}}

\begin{document}
\maketitle

%------------------------------------------------
\section{Motivation}

\begin{frame}{Why Not Fermat?}

\begin{itemize}
    \item Fermat test is fast but unreliable
    \item Carmichael numbers fool it for all bases
\end{itemize}

\vspace{1em}

\begin{block}{Goal}
Design a test that:
\begin{itemize}
    \item Remains fast
    \item Dramatically reduces error
\end{itemize}
\end{block}

\end{frame}

%------------------------------------------------
\section{Key Idea}

\begin{frame}{Structure of $n-1$}

Let $n$ be odd. Write:
\[
n-1 = 2^s d \quad \text{with } d \text{ odd}
\]

\vspace{1em}

\begin{itemize}
    \item This decomposition is unique
    \item Central to Miller--Rabin test
\end{itemize}

\end{frame}

%------------------------------------------------

\begin{frame}{Core Observation}

If $n$ is prime and $\gcd(a,n)=1$:
\[
a^{n-1} \equiv 1 \pmod{n}
\]

\vspace{1em}

So:
\[
a^{2^s d} \equiv 1 \pmod{n}
\]

\end{frame}

%------------------------------------------------

\begin{frame}{Chain of Squares}

Define:
\[
x_0 = a^d, \quad x_1 = a^{2d}, \quad x_2 = a^{4d}, \dots
\]

\[
x_{i+1} = x_i^2 \pmod{n}
\]

\vspace{1em}

Eventually:
\[
x_s = a^{2^s d} \equiv 1 \pmod{n}
\]

\end{frame}

%------------------------------------------------

\begin{frame}{Critical Property (Primes)}

If $n$ is prime:

\begin{itemize}
    \item Either $x_0 \equiv 1 \pmod{n}$
    \item Or some $x_i \equiv -1 \pmod{n}$
\end{itemize}

\vspace{1em}

\begin{block}{Why?}
Because $-1$ is the only nontrivial square root of $1$ modulo a prime
\end{block}

\end{frame}

%------------------------------------------------
\section{Key Lemma}

\begin{frame}{Nontrivial Square Roots}

\begin{block}{Lemma}
If $p$ is prime and:
\[
x^2 \equiv 1 \pmod{p}
\]
then:
\[
x \equiv \pm 1 \pmod{p}
\]
\end{block}

\end{frame}

%------------------------------------------------

\begin{frame}{Proof}

\[
x^2 \equiv 1 \Rightarrow (x-1)(x+1) \equiv 0 \pmod{p}
\]

\vspace{1em}

Since $p$ is prime:
\[
p \mid (x-1) \quad \text{or} \quad p \mid (x+1)
\]

\vspace{1em}

\[
\Rightarrow x \equiv 1 \ \text{or} \ -1 \pmod{p}
\]

\end{frame}

%------------------------------------------------

\begin{frame}{Implication}

\begin{itemize}
    \item If we ever find:
    \[
    x_i^2 \equiv 1 \pmod{n}
    \]
    but
    \[
    x_i \not\equiv \pm 1
    \]
\end{itemize}

\vspace{1em}

\begin{block}{Conclusion}
$n$ must be composite
\end{block}

\end{frame}

%------------------------------------------------
\section{Miller--Rabin Test}

\begin{frame}{Algorithm}

\begin{enumerate}
    \item Write $n-1 = 2^s d$
    \item Pick random $a$
    \item Compute $x = a^d \bmod n$
    \item If $x=1$ or $x=n-1$, continue
    \item Repeat squaring:
    \[
    x \leftarrow x^2 \bmod n
    \]
    \item If $x = n-1$, continue
    \item If $x = 1$, return composite
    \item Otherwise, return composite
\end{enumerate}

\end{frame}

%------------------------------------------------

\begin{frame}{Example: $n = 561$ (Carmichael Number)}

\[
561 = 3 \cdot 11 \cdot 17
\]

\[
n - 1 = 560 = 2^4 \cdot 35
\]

Choose $a = 2$

\vspace{0.5em}

\textbf{Step 1:}
\[
x_0 = 2^{35} \bmod 561 = 263
\]

\vspace{0.5em}



\end{frame}
\begin{frame}
\textbf{Step 2: Repeated squaring}
\[
\begin{aligned}
    x_1 &= 263^2 \bmod 561 = 166 \\
    x_2 &= 166^2 \bmod 561 = 67 \\
    x_3 &= 67^2 \bmod 561 = 1
\end{aligned}
\]
\begin{block}{Observation}
We found $x_i = 1$ but previous value $\neq \pm1$
\end{block}

\[
67 \not\equiv \pm 1 \pmod{561}
\]

\vspace{0.5em}

\begin{block}{Conclusion}
561 is \bld{composite}
\end{block}

\end{frame}

%------------------------------------------------

\begin{frame}{Witnesses}

\begin{itemize}
    \item If test declares composite $\Rightarrow$ correct
    \item Such $a$ is called a \itl{witness}
\end{itemize}

\vspace{1em}

\begin{block}{If no witness found}
$n$ is \bld{probably prime}
\end{block}

\end{frame}

%------------------------------------------------
\section{Correctness}

\begin{frame}{Error Probability}

\begin{block}{Theorem}
If $n$ is composite, then at least $\frac{3}{4}$ of choices of $a$ detect it
\end{block}

\vspace{1em}

\[
\Rightarrow \text{Error probability} \le \left(\frac{1}{4}\right)^k
\]

After $k$ independent iterations:
\[
\Pr[\text{error}] \le \left(\frac{1}{4}\right)^k
\]

\vspace{0.5em}

For $k = 5$:
\[
\Pr[\text{error}] \le \left(\frac{1}{4}\right)^5 = \frac{1}{1024} \approx 0.1\%
\]
\end{frame}

%------------------------------------------------

\begin{frame}{Comparison with Fermat}

\begin{itemize}
    \item Fermat can fail for all $a$ (Carmichael)
    \item Miller--Rabin fails for at most $1/4$ of $a$
\end{itemize}

\vspace{1em}

\begin{block}{Conclusion}
Exponentially stronger reliability
\end{block}

\end{frame}

%------------------------------------------------
\section{Complexity}

\begin{frame}{Time Complexity}

\begin{itemize}
    \item Modular exponentiation: $\bigO((\log n)^3)$
    \item $k$ iterations:
\end{itemize}

\[
\bigO(k (\log n)^3)
\]

\vspace{1em}

\begin{block}{Practical}
Very fast and widely used
\end{block}

\end{frame}

%------------------------------------------------

\begin{frame}{Takeaway}

\begin{itemize}
    \item Fermat: fast but unreliable
    \item Miller--Rabin: fast and highly reliable
\end{itemize}

\vspace{1em}

\begin{block}{Big Picture}
Randomization + structure $\Rightarrow$ powerful primality testing
\end{block}

\end{frame}

\end{document}